\begin{minipage}[t][][t]{0.60\linewidth}
    On réalise l'expérience des rails de Laplace décrite dans le cours. Le
    champs magnétique $\vec{B}$ est dirigé vers le haut. Le tige métallique AB
    peut glisser sans frottement en restant perpendiculaire aux rails.
\end{minipage}
\hfill
\begin{minipage}[t][][t]{0.38\linewidth}
    \vspace{-5mm}
    \begin{figure}[H]
        \centering
        \begin{tikzpicture}[]
            \draw[] (0,0) coordinate(A) -- ++(-5,0) coordinate(B)
                to[battery2] ++(0,-2) coordinate(C) -- ++(5,0) coordinate(D);
            \draw[line width=2pt,shorten <=-2pt,shorten >=-2pt]
                ($(A)!.2!(B)$) node[above] {$A$} --
                ($(D)!.2!(C)$) node[below] {$B$};
            \node[circle,draw,inner sep=2pt,thick,xshift=3cm,
                label=$\vec{B}$] (O) at
                ($(B)!.5!(C)$) {};
            \node[circle,fill,inner sep=0.8pt] at (O.center) {};
        \end{tikzpicture}
    \end{figure}
\end{minipage}

\vspace{-5mm}
\begin{enumerate}
    \item Indiquer dans quel sens le courant circule dans la
          tige.~\textbf{(0,5pt)}
    \item Indiquer le sens de la force de Laplace (sens du mouvement de la barre
          sur les rails).~\textbf{(0,5pt)}
    \item Calculer la valeur de la force de Laplace avec~:~\textbf{(1pt)}
          \[
              I=\qty{1.5}{\A}~;~B=\qty{85}{\mT}~;~\ell=\qty{10}{\cm}.
          \]
\end{enumerate}


